\subsection{The full topology of framed squares}

\begin{result}{E.40}{Proved}{Tripod squares recover every full framed transporter}
Attach a rational tripod to each labelled boundary of a stable
bordered core, giving its two other punctures new distinct labels.
Write $t$ for the resulting tag multicurve.  For essential interior
core multicurves $\gamma_0\supseteq\gamma_1$, write
$D=\gamma_0\setminus\gamma_1$ and put
\[
 A_i=t\cup\gamma_i,\qquad B_i=\gamma_i,
 \qquad \eta_i:A_i\longrightarrow B_i.                       \tag{E.160}
\]
Here $\eta_i$ smooths exactly the tag nodes.  In the strict completed
exit category, the space of commutative squares from $\eta_0$ to
$\eta_1$, with the indicated horizontal incidence, is canonically
the full completion of the bordered-core transporter
\[
 \operatorname{Trans}_{\Stab(t)}(\gamma_0,\gamma_1)
                         /\Delta_{\gamma_0}.                 \tag{E.161}
\]
For one orbit it is the compact homogeneous space
\[
 (\widehat C_1\times\widehat C_0)/\widehat S,
 \quad C_i=\Stab(t,\gamma_i)/\Delta_{\gamma_i},\quad
 S=\Stab(t,\gamma_0,\gamma_1)/\Delta_{\gamma_1}.               \tag{E.162}
\]
All incidence and pinning choices give the corresponding finite
disjoint union.  The comparison identifies both horizontal
projections, all endpoint actions and compositions, and the
projections of specified successive smoothing flags.  In particular discrete framed squares
are dense in this actual completed square space.

The annular unit is represented by the single four-point crossing
of E.35; its two boundary twists are the same positive meridian.
The assertion for stable cores uses their distinct tag curves as
in (E.160).
\end{result}
\begin{proof}
We first give the discrete comparison, including its inverse.
Every tag encloses its own specified pair of new labels.  An
interior essential core curve cannot enclose just this pair:
before attachment it would be parallel to the corresponding core
boundary.  Nor can any old rational-tail curve have this pair of
new labels.  Hence a transporter $A_0\to A_1$ preserves every tag
individually, even when it permutes other curves or cut components.
It has a representative $c\in\Stab(t)$.

The two composites around a square agree modulo
$\Delta_t\Delta_{\gamma_0}$ in the transporter from $A_0$ to
$B_1$.  Thus a representative of its bottom arrow differs from
$c$ by an element of this lattice and also lies in $\Stab(t)$.
The bottom arrow determines the top one by the additional quotient
by $\Delta_t$.  Explicitly the comparison is
\[
 c\Delta_{\gamma_0}\longmapsto
       (c\Delta_{t\cup\gamma_0},c\Delta_{\gamma_0}).          \tag{E.163}
\]
Its inverse is the bottom representative, taken in the framed
transporter.  Changes of representative are exactly
$\Delta_{\gamma_0}$, so both maps are well-defined and inverse.
Extension by the standard tripods and restriction to the core
identify $\Stab(t)$ with the boundary-fixed core group, by the
literal cutting presentation (E.84).  This proves (E.161) on the
discrete category, rather than only on its isotropy groups.

Set $G_i=\Aut(B_i)$ and $H_i=\Aut(A_i)$.  Normalize at
$\gamma_i$, retaining the chosen branch pinnings.  On each
normalized component, $\eta_i$ is a crossing smoothing a specified
collection of rational tails.  Components without a tag have the
identity crossing.  E.33, applied component by component, gives
\[
 \Aut(\eta_i)=\widehat C_i
       \lhook\joinrel\longrightarrow\widehat G_i\times
                                      \widehat H_i.           \tag{E.164}
\]
Its first projection is injective; its second is onto, with kernel
$\widehat\Delta_t$.  To retain unpinned component or branch
permutations, first make all the finitely many pinnings and then
induce over them.  The kernel calculations are unchanged.  This
is finite-index passage on the component groups and their
morphism torsors, so their full completions, not smaller
congruence quotients, occur in (E.164).

We need the entire composition fibre, not merely (E.164).
Let $z_0:A_0\to B_1$ be the composite smoothing $t\cup D$.
Cut at the retained multicurve $\gamma_1$.  A choice of its
persistent branches gives precisely the exit category of its
normalized components, with $B_1$ as its open object and residual
source multicurve $t\cup D$.  We first justify the finite-pinning completion used here.  On
the finite orbit skeleton let $P(i)$ be the finite set of
submulticurves of its representative $\sigma_i$ with the
specified retained type, together with their finitely many
branch pinnings.  An arrow $g:i\to j$ pulls a target marking
$\rho$ back to $g^{-1}\rho$.  This is independent of the
source-twist representative, so $P$ is an actual finite
contravariant functor.  Its category of elements has the
marked objects and a unique lift of $g$ for each target
marking.  Given a finite categorical quotient $F$ of this
category of elements, record for every downstairs arrow $g$
its map $P(g)$ and the finite array of all values
$F(g_\rho)$, $\rho\in P(j)$.  The array for $hg$ has
entries
$F(h_\rho)F(g_{P(h)\rho})$.
Equality of arrays is therefore a finite categorical
congruence downstairs.  Its pullback refines $F$.
Conversely a downstairs finite congruence pulls back after
retaining the finite $P$ coordinates.  Thus the two systems
of finite quotients are cofinal.  Completion commutes with
this specified finite marking construction, including all
its projection maps.  Equivalently each fixed-lift Hom-space
is the indicated clopen fibre of $P(g)$ in the completed
downstairs Hom-space; discrete lifts are dense in it.

In this marked category the comparison cuts each
representative along those persistent curves; its inverse extends
the component representatives over the same pinned collars.
The twists of those collars are killed on both sides.  These
operations are inverse on arrows and their compositions, as in
A.9--A.10 and the admissible square of A.15.  Keeping the
persistent marking gives a convex component diagram: an
intermediate contraction cannot erase a node which its final
target retains.  The convex finite-quotient
extension argument (E.86) applies to this marked component
diagram.  The array construction just proved then retains
all its finite quotients in the unpinned completed category,
with the finite permutation actions unchanged.

Apply E.29 in that component diagram.  On the chosen incidence
orbit it gives
\[
 \Hom(A_0,B_1)=E/\widehat K,
 \quad E=\widehat G_1\times\widehat H_0,
 \quad K=\Stab(t,\gamma_0,\gamma_1)/\Delta_{\gamma_1}.        \tag{E.165}
\]
The equality in (E.165) denotes that specified orbit; other
matching contraction flags retain their own finite labels.
The two maps from $K$ to the endpoints are its target inclusion
and source quotient, exactly as in (E.80).

Inside this composite, the intermediate multicurve $B_0$ is
obtained by dropping all the distinguished tag curves; $A_1$ is
obtained by dropping the remaining residual set $D$.
Both intermediate choices are therefore intrinsic.  The fresh
labels rule out an exchange of any curve in $D$ with a tag,
including when such a curve bounds a rational tail with old
labels.  Curves within $D$ may be permuted; this does not
change either of the two intermediate subsets.  Consequently
the automorphism group of either specified two-step
factorization projects isomorphically to $K$.  Its intermediate
coordinate is the very same cut mapping class modulo the
appropriate intermediate twist lattice.  This explicitly gives
the inverse to that projection.

Here is a direct completion check for these coarse two-step
factorizations, including arbitrary $D$.  The stabilizer of
$B_0\to B_1$ has source map to $G_0$ with finite-index image
and kernel
$\Delta_{\gamma_0}/\Delta_{\gamma_1}\cong\mathbb Z^D$.
This kernel has its full topology in the stabilizer, because
E.29 gives its full topology in the normalized target component
group.  The stabilizer of the two-step factorization through
$B_0$ is the fibre product of this arrow stabilizer and $C_0$
over $G_0$.  Apply the general fibre-product lemma (E.87),
with its finite-index variant.  It identifies the actual
completed fibre product with the full completion of the
discrete two-step stabilizer, which the preceding paragraph
identified with $K$.  All the coarse individual arrow
stabilizers used here are obtained by E.29 after normalizing
at their retained target multicurve, with the finite-pinning
comparison proved above.  The proof of (E.87) applies to any
kernel having full topology, so it applies to the entire
lattice $\mathbb Z^D$.

For the factorization through $A_1$, form the fibre product
of the stabilizer of $A_0\to A_1$ and $C_1$ over $H_1$.
Use the surjective source map $C_1\to H_1$; its kernel
$\Delta_t$ has full topology by E.33.  Equation (E.87) again
identifies this stabilizer with the full completion of its
discrete group $K$.  These are the natural projections of
tuples, so both maps to the total-composite stabilizer in
(E.165) are the completed discrete isomorphism onto $K$.
No ordering of the curves of $D$, or restriction to their
pointwise stabilizer, has been made.

For compatibility with refinements, refine a tag-smoothing
arrow in any fixed order.  By (E.164), an automorphism of
that coarse arrow is an element of $\widehat C_i$.  Its
actual quotient coordinates on the intermediate subsets,
given by E.33's Boolean tail category, extend it uniquely
to the whole flag.  Uniqueness follows from E.105 and the
injective target projection in (E.164).  Flags of core
smoothings are computed by successive applications of
(E.87), or by E.32 for an adjacent refinement.  Their finite
order and pinning choices are retained, including the
permutations of curves in $D$.  This proves the needed
coarse and refined flag comparisons without an interchange
assertion for arbitrary diagram categories.

We next check that there are no additional completed
factorization orbits.  For any fixed finite object list the
discrete composable-arrow tuples have finitely many orbits
under the endpoint and intermediate automorphism groups.
They are the actual finite labelled matching flags of A.13.
Each corresponding profinite group orbit is compact, hence
closed in the finite product of completed Hom-spaces.
Their finite union contains the product of the dense
discrete Hom-spaces, so it is the entire product.  The preceding fibre-product calculation
computes the full stabilizer of each coarse two-step flag;
its refinements retain the same actual projection maps.  Different composite-incidence
labels remain disjoint by (E.165).  Within the chosen
composite orbit the two intrinsic intermediate choices
therefore give exactly the two factorization spaces
\[
 \begin{split}
 \mathcal F_0&=(E\times\widehat G_0)\big/
        \operatorname{Graph}(\widehat K\to\widehat G_0),\\
 \mathcal F_1&=(E\times\widehat H_1)\big/
        \operatorname{Graph}(\widehat K\to\widehat H_1).
 \end{split}                                                  \tag{E.166}
\]
Each graph here also includes the known embedding
$\widehat K\to E$.  Composition is the map induced by
projection to $E/\widehat K$.  The formula proves, in
particular, that the fibre over a fixed composite is the
indicated middle endpoint group.  We used density in an
ordinary product of Hom-spaces; no density assertion about
an unspecified equalizer was used.

Now take an actual completed square with common composite
$z$.  Choose $(g_1,h_0)\in E$ carrying $z_0$ to $z$.
Let $\alpha:B_0\to B_1$ and
$\alpha^{\rm top}:A_0\to A_1$ be the two pinned horizontal
arrows of the canonical square.  Formula (E.166) supplies
middle coordinates $x\in\widehat G_0$ and
$y\in\widehat H_1$ for its two factorizations, whose
ordered arrows are
\[
 \begin{split}
 &(x\eta_0h_0^{-1},\ g_1\alpha x^{-1}),\\
 &(y\alpha^{\rm top}h_0^{-1},\ g_1\eta_1y^{-1}).
 \end{split}                                                  \tag{E.167}
\]
Their vertical arrows are the fixed $\eta_0,\eta_1$ exactly
when
\[
 (x,h_0)\in\widehat C_0,
          \qquad (g_1,y)\in\widehat C_1.                     \tag{E.168}
\]
This is (E.164), including both endpoint actions.
Consequently every completed square is the endpoint
translate of the canonical square by an element of
$\widehat C_1\times\widehat C_0$.

The stabilizer of that square is the completion of $S$ in
(E.162), embedded by its actual intermediate quotient
coordinates.  One can see this directly in (E.166): its
stabilizer is the same $\widehat K$, now with the two
coordinates forced by (E.168); discretely $K=S$.
Its full topology in $C_1\times C_0$ also has a separate
check.  The map $S\to C_0$ has finite-index image and
kernel the residual twist lattice
$\Delta_{\gamma_0}/\Delta_{\gamma_1}\cong\mathbb Z^D$.
That lattice has its full topology in $C_1$: after
normalization at $\gamma_1$ its curves are essential and
nonparallel on the cut components, and E.29 supplies all
of its finite quotients.  Its topology in the framed
subgroup is at least this full topology.  The diagonal argument of
E.32 therefore embeds $\widehat S$ into
$\widehat C_1\times\widehat C_0$.  The compact-orbit map
now identifies the square space homeomorphically with
(E.162); its inverse is given by the coordinates
(E.167)--(E.168) modulo $\widehat S$.

Finally this homogeneous space is precisely the full
categorical completion of (E.161).  In the bordered-core
category retain the target multicurve $\gamma_1$ by the
same finite-pinning construction and cut along it.  Its
chosen incidence is a crossing to the open object of the
cut-component diagram.  The direct proof of E.29 applies
also to these boundary-fixed groups: their independent
multitwists have full topology by extension into their
labelled tripod surface, and every finite group quotient
produces exactly the finite transporter category (E.82).
The diagonal topology and finite coset argument of
(E.80)--(E.83) therefore identify this incidence completion
with (E.162).  The array congruences above retain all these
quotients in the whole framed category.  This proves the
claim for every finite $D$; for an adjacent pair the
convex-pair argument (E.86)--(E.89) is the same computation.
The image of $C_1\times C_0$ is dense in this compact
homogeneous space, proving the density assertion for
squares as well.

All formulas used the actual representatives and their
quotients.  Projection to either horizontal arrow in
(E.163) and (E.167) is therefore the actual projection of
the square.  Refining a succession of horizontal arrows
uses the corresponding actual matching flags; E.32 and
the unique Boolean tag extension above identify their
stabilizers and all their projections.  Discrete composition
is the identical multiplication of core representatives in
(E.163).  Products of discrete Hom-spaces are dense, so the
continuous comparisons commute with every completed
composition as well.  This proves the full framed-category
comparison, its endpoint actions and its flag compatibility,
without an interchange assertion for arbitrary categorical
completions and arrow categories.
\end{proof}
